\section{Limitations and Conclusion}
\label{sec:limitations}

This study is predictive, not causal. The graph model shows that observed
proximity, combat, and co-pressure relations help predict next-phase
elimination hazard; it does not prove that changing a team's strategy to
alter one relation would causally change survival. The next-phase prediction
task is also an operational evaluation of the representation, not the final
purpose of the paper.

Four limitations follow. First, telemetry contains only observable positions,
states, and events. It does not directly observe terrain, indoor/outdoor
status, line of sight, voice communication, player skill, or intention. Second,
team centroids and phase graphs compress complex squad behavior. Team spread
and nearest-enemy features partially compensate, but the same centroid can
hide different formations and cover conditions. Third, the population depends
on telemetry quality and graph eligibility. Most excluded matches are tutorial
or training-range sessions, but missing position traces and too-short phase
sequences still remove a small number of records. The random match-level split
tests generalization inside the observed distribution; it does not guarantee
robustness across future patches, seasons, regions, or tournament settings.
Fourth, later phases contain fewer teams and target cells, increasing
variance and survivor bias, and graph construction uses fixed distance and
pressure thresholds rather than a full sensitivity sweep.

We reconstructed PUBG telemetry as a phase-aligned STDB and used it to define
a team-level next-phase elimination hazard task for evaluation. The point of
the prediction task is not merely to guess which team will be eliminated; it is
to test whether a representation that preserves movement, zone contraction,
combat events, and team relations captures the risk state of a battle royale
match. On 21,824 matches, the phase-level STDB representation carries more
risk information than compact strategy-score summaries, and dynamic team
graphs add a smaller but statistically reliable relational signal beyond
feature-only models. The result supports treating PUBG telemetry as a
spatio-temporal survival process rather than only a source of post-match
aggregate statistics.
